\documentclass{article}
\usepackage{amsmath,amssymb,hyperref}
\title{Singleton module actions: a cone-preserving enlargement}
\date{15 November 2026}
\begin{document}\maketitle
\paragraph{Status.} Written argument independently reviewed in the stated scope; not a mechanized proof or a full completion claim. This route replaces indiscriminate free adjunction admission, which can destroy old synchronization limits. It does not yet extend directed parallel or prove the full relative feedback construction.

\paragraph{Independent hypotheses.} Let $Q$ be a suitably universe-bounded locally cocomplete bicategory with separately cocontinuous composition. Assume the small enriched module bicategory $\mathcal M=\operatorname{Mod}(Q)$ exists, with the indicated hom categories complete enough that postcomposition by any module has a right adjoint. The locally presentable matrix-action argument below verifies the required column residuals for the actual $PK$ base. Let $*_x$ denote the singleton $Q$-category at extent $x$.

\paragraph{Discharging residual existence for the actual local presheaf base.} For $Q=PK$, fix a small enriched category $A$ and extent $x$. Columns are algebras for the matrix-action monad on
\[
 \prod_{a\in\operatorname{ob}A}Q(x,|a|),\qquad
 (S W)_b=\coprod_a A(b,a)\circ W_a.
\]
The monad unit and multiplication are the enriched identity and composition maps. Every hom of $PK$ is a presheaf category. After choosing a common regular cardinal above the small indexing data, the product is locally presentable and the displayed monad is accessible. In fact it preserves colimits, by separate cocontinuity of composition. Its algebra category is therefore locally presentable; the forgetful functor creates the colimits in question. Postcomposition by a module is constructed from coproducts and the module-action coequalizer, componentwise. It consequently preserves colimits and has a right adjoint by the locally presentable adjoint functor theorem. This verifies the hypothesis here without incorrectly inferring all needed ends from biclosedness alone.

\paragraph{The action representation.} Define
\[
 R(A)(x)=\mathcal M(*_x,A),\qquad R(M)_x=M\odot -.
\]
An arrow $u:x\to y$ of $Q$ acts by precomposition, from $R(A)(y)$ to $R(A)(x)$. Thus $R(A)$ is a pseudofunctor $Q^{op}\to\mathrm{Cat}$, reversing one-cells but not two-cells. Associativity makes $R(M)$ pseudonatural. Take as enlarged target the bicategory whose objects are these pseudofunctors and whose arrows are lax transformations, retaining modifications. Extra objects may be admitted later; the present recovery assertion concerns the tagged images $R(A)$.

\paragraph{Native recovery, not defeated by extra arrows.} A modification between $R(M)$ and $R(N)$ can be evaluated at the representable column $A(a,-):*_x\to A$. The module Yoneda isomorphisms give a column map $M(-,a)\to N(-,a)$. Naturality under the $Q$-action and the enriched category actions supplies exactly a module map. Conversely every module map supplies a modification. To show these are inverse on arbitrary columns, use the enriched density coequalizer expressing a column as coproducts of representable columns composed with its component arrows. Both $M\odot-$ and $N\odot-$ preserve these colimits; the modification is determined on the generators by ordinary naturality. For action compatibility explicitly apply pseudonaturality to the arrow $A(a',a)$ and naturality to the canonical action map from the corresponding representable composite. The left target action already holds since the evaluated cells are column-module maps. Adding lax transformations as new one-cells does not itself add modifications between the old pseudonatural transformations: full faithfulness of a hom functor allows new objects outside its image.

\paragraph{Adjoints and variance.} Write a lax structure cell for $\alpha:F\to G$ along an arrow $v:i\to j$ in the indexing bicategory as
\[
 G(v)\alpha_i\Longrightarrow\alpha_jF(v).
\]
If $\alpha$ is pseudonatural and each $\alpha_i\dashv\beta_i$, the mate of the inverse structural isomorphism is
\[
 F(v)\beta_i\Longrightarrow\beta_jG(v).
\]
These cells make $\beta:G\to F$ lax. The mate correspondence takes the pseudonaturality unit/composition equations to the lax equations, and makes the component units and counits modifications. Hence $\alpha\dashv\beta$ in the lax-transformation bicategory. In particular, every represented module arrow has such a right adjoint under the stated residual hypothesis. The right adjoint need not be a represented module or pseudonatural. This is doctrinal adjunction, not a new adjoint-existence theorem.

\paragraph{Synchronization preservation.} Let $D\to\mathcal M$ be a pseudodiagram with an incoming lax bilimit $E$. Each represented hom $\mathcal M(*_x,-)$ preserves its cone equivalence, so $R(E)(x)$ is the pointwise lax limit of the corresponding categories. This is still a lax bilimit in the enlarged target, tested against arbitrary lax transformations, not merely the old module maps. Construction: transport each $Q$-action across the pointwise cone equivalence using the pseudonaturality of the diagram arrows. An incoming lax transformation to this pointwise object is a family of functors to the pointwise limits with structural natural transformations; the universal property identifies these with compatible families of lax transformations to the diagram. Modifications are identified componentwise. The action laws and structural equations can be checked after every cone projection, which is jointly faithful on the relevant natural transformations. This argument requires the diagram arrows to be pseudonatural; it does not assert unrestricted pointwise creation for arbitrary diagrams of lax arrows.

\paragraph{What this supplies.} The canonical output comparison in the earlier synchronization theorem is an old module arrow. Its represented image acquires a right adjoint without losing the specific source and target cone universal properties. This avoids the defect of the free compact envelope. The resulting synchronization projections may still be noninvertible native cells. Directed parallel, assembly Beck--Chevalley and preservation of the tagged pure interface are further obligations, not consequences of the word representation. In particular a generic right adjoint is not cocontinuous, so the previously established coend extension of parallel cannot simply be reused on all new lax arrows.

\paragraph{A parallel-extension shortcut fails its unit test.} This is a diagnostic, not a no-go theorem. In the ordinary presheaf realization, an arbitrary functor $F:\widehat A\to\widehat B$ has a cocontinuous replacement $C(F)=\operatorname{Lan}_{y_A}(F y_A)$ with a canonical map $C(F)\to F$. Applying the old module tensor only after this replacement gives $C(F)$, not $F$, upon tensoring with the old tensor unit. Already for $A=B=1$, the right adjoint $F(X)=X^{\{0,1\}}$ satisfies $F(1)=1$, so $C(F)(X)=X$, whereas $F(2)$ has four elements. Thus this replacement cannot furnish an invertibly unital extension to all new arrows. Using Cartesian product in Cat instead changes the old tensor unit from Set to the terminal category and is not the required recovery either. A genuine extension needs additional structure or a different construction; neither failure excludes one.

\paragraph{Source trail.} The ordinary $\mathrm{Prof}\to\mathrm{Cat}$ action representation and enriched density are established presheaf machinery. The relevant higher adjunction principle is also standard; a modern primary treatment is \emph{Straightening for lax transformations and adjunctions of $(\infty,2)$-categories}, \url{https://doi.org/10.1007/s00029-025-01084-z}. Its theorem must be compared with the explicit convention above rather than used to infer our many-object recovery and limit assertions. The prior local/module construction uses Kelly--Labella--Schmitt--Street, \emph{Categories enriched on two sides}, and Garner--Shulman, \url{https://arxiv.org/abs/1301.3191}. No novelty claim is made.
\end{document}
